\chapter{ADS-24 Data Construction}
\label{app:ads24-data}
% 개요 15장(부록 D). label app:ads24-data는 08_ads24.tex의 참조에 맞췄다.
% 근거: RECORD:236--254 (자료와 과제 구성의 세부), RECORD:265--275 (외부 자료 선정 경위).
% 완화한 세 조건(발생 시각 순서, 자료의 단계 라벨, 사용자 조건 보류)은 08장 첫 문단에
% 있으므로 여기서 되풀이하지 않았다.

Table~\ref{tab:ads24-construction} traces the ADS-24
data~\cite{nyberg2026cyber} from the episodes in the official summary to the
132 runs of the contract. The counts in the table are our own tallies
from the downloaded files. One episode, identified by its episode
identifier, is one run; we use \emph{run} throughout. Each episode ships
with an attack log, called Lore in the dataset, that lists the
attacker's stages with their tactics; 255 runs have one. A run became
usable for the task only after its stages were placed on the time axis
of its Wazuh alerts, and this step dropped 288 stages that fell outside
the extraction window and one run whose Wazuh JSON could not be parsed.

The stored preflight report documents how stages and alerts are placed
on one time axis. Lore stage start times carry no time zone and are
read as Europe/Stockholm local time, which the dataset's own utility
code declares; Wazuh alerts are read by their UTC timestamp. The cutoff
of a boundary is the Lore start time of the next tactic segment, and
only alerts timestamped strictly before it enter the input, in the
original Wazuh order. When several stages start in the same second,
the one with the smallest action identifier is kept, and a stage with
several tactics takes the first. Stages are not paired one-to-one with
alerts. As a check of the common time base, the report compares, in 68
episodes, the Lore time of a CrackMapExec action with the Wazuh time of
the matching TCP/445 event on the defended host; the Wazuh time was
later by a median of 11.47~s. We did not rerun this procedure, because
its code and its per-run audit file are not in the archive available to
us, and no binding was checked by hand. Whether the Wazuh timestamp is
a receipt time is not established, and a different binding could move
segment boundaries and next-tactic labels; this is one reason
Chapter~\ref{ch:ads24} is an exploratory check.

The contract population then excludes runs flagged with
\texttt{malicious\_users} or \texttt{block\_blinds\_blue} and runs whose
defender is unknown. These three flags leave 134 runs, the same number
as the 134 episodes on which the dataset paper bases its
results~\cite{nyberg2026cyber}. We did not find a list of episode
identifiers in that paper, so the two sets were matched in number only. Two of the 134 runs are then dropped because one of their two
tactic stages falls after the end of the Wazuh extraction, which leaves
the 132 runs on which every ADS-24 result in this thesis is computed.
Two label baselines, computed on the 132 runs with five folds, show how
much the labels alone reveal. They are not the FIT prior against which
Chapter~\ref{ch:ads24} compares the models
(Table~\ref{tab:ads24-groups}).
% 출처: sources/atlas_answers/ads.md 25·26번.
%   atlas OT/PROPHET_ENV_CONSENSUS_KD_V3_20261120/ADS24_V3_DRAFT.md:7 (독립 단위 episode_id);
%   ADS24_PREFLIGHT.md:19 (Europe/Stockholm, +02:00 254개), :23 (60초 이내 SMB 흔적, 일대일 결속 아님),
%   :29-31 (제외 목록), :35 (같은 time_start는 최소 action_id), :76 (receipt 필드 미확인),
%   :80 ("논문의 134와 숫자만 일치", 2개 빠져 132); ADS24_CONTRACT_DRAFT.md:27 (cutoff와 같은 시각 제외).
%   68 에피소드·중앙값 11.47초: RECORD(inputs/drop)의 시각 앵커 값, 정의는 PREFLIGHT:23.
% 8회차 인용 P1-2: 134 episode는 nyberg2026cyber 본문에서 확인됐다(인용 리뷰). 134 바로 뒤에 인용을 달고,
%   식별자 목록이 없다는 말은 "우리가 찾지 못했다"로 한정했다. 해결된 TODO는 지웠다.

Only alert identity, severity, host role, and timing enter the input;
every field that names the answer or the run was removed. In particular, Wazuh rules can map an alert to
ATT\&CK techniques and tactics through the \texttt{rule.mitre.*}
fields,\footnote{Wazuh documentation on the MITRE ATT\&CK mapping of
rules: \url{https://documentation.wazuh.com/current/user-manual/ruleset/mitre.html}.}
which would name the answer. In our parse of
the data these fields are present on 1,252,951 alerts, and they were
removed.

Before any model was scored, we fixed the segment rule of
Section~\ref{sec:ads24-task} because of how often a stage repeats its
predecessor: if stages are not
merged, about 62.7\% of consecutive boundaries carry the same tactic as
the stage before them, and a model could score well by copying the last
label. The alert window of 512 and the 1,346 updates were chosen by an
inner validation inside FIT, with seeds outside the contract, before
stage~1.
% 근거: RECORD:238--252, 254. 수치는 기록의 값을 그대로 옮기고, 기준선 NLL은 본문 규칙대로
% 셋째 자리로 줄였다(1.8547→1.855, 1.4965→1.497).
% 6회차 문장 P2-4: oracle boundary 정의는 8.1절에만 둔다.

\begin{table}[htbp]
\centering
\small
\caption{Construction of the ADS-24 task, from the official episode
summary to the inputs of the contract. Run and boundary counts per role
are those of Section~\ref{sec:ads24-task}. Label baselines are computed
on the 132 contract runs with five folds and differ from the FIT-only
baseline of Table~\ref{tab:ads24-groups}.}
\label{tab:ads24-construction}
\begin{tabularx}{\textwidth}{@{}>{\raggedright\arraybackslash}p{0.30\textwidth}>{\raggedright\arraybackslash}X@{}}
\toprule
Item & Value \\
\midrule
Episodes in the official summary & 300 \\
Runs with an attack log (Lore) & 255, of which 254 have at least two
  tactic stages \\
Runs usable after time binding & 251 (288 stages outside the extraction
  window and one run with unparsable Wazuh JSON excluded) \\
Parsed Wazuh alerts & 2,067,288 \\
Time offset & UTC+02:00 confirmed in 254 runs; the whole collection
  period lies within daylight saving time \\
Time base & Lore start times as Europe/Stockholm local time; Wazuh UTC
  timestamps; alerts strictly before the cutoff. Check: Wazuh TCP/445
  time minus Lore CrackMapExec time, 68 episodes, median 11.47~s; not
  rerun by us \\
Contract population & 134 runs after the flags \texttt{malicious\_users},
  \texttt{block\_blinds\_blue}, and unknown defender; 132 after two runs
  whose second stage falls after the Wazuh extraction end \\
Runs / boundaries by role & FIT, no users: 48 / 8,461; preservation, no
  users: 15 / 2,982; unlabeled adaptation, with users: 34 / 5,811;
  environment evaluation, with users: 35 / 5,639 \\
Label baselines (NLL / accuracy) & majority class 1.855 / 32.0\%;
  previous-tactic bigram 1.497 / 42.2\% \\
Boundaries with no new alert & 14.7\% (over the 251 usable runs) \\
Removed from the input & \texttt{rule.mitre.*} fields (present on
  1,252,951 alerts); 227 alerts from the flightlogs agent; free text;
  run, scenario, defender, and user identifiers; \texttt{rule.groups}
  (empty allow list, so all zero) \\
Input fields & \texttt{rule.id} (FIT vocabulary), \texttt{rule.level},
  host role, interval between alerts, age at the cutoff, alert-presence
  indicator \\
Inner selection & FIT 40/8 inner validation over 9 candidate pairs,
  seeds outside the contract: window 512, 1,346 updates \\
\bottomrule
\end{tabularx}
\end{table}
% 근거: RECORD:240--252. 역할별 수는 08장과 같은 값이다(ADS24_STAGE2_REPORT.md "계보").

Among the seven public datasets we examined, none allowed the PROPHET
task to be posed as it stands, so the search admitted datasets under the
three relaxations stated at the start of Chapter~\ref{ch:ads24}. Table~\ref{tab:ads24-candidates} lists
the seven candidates that were examined. The deciding requirements were
practical ones: independent runs with stage-level labels and times,
benign user activity during the attack, a license and download path we
could confirm, and a held-out condition that does not leave most test
labels unseen in training. CAM-LDS illustrates the last point. It is a
Linux collection of 34 runs without benign user activity during the
attacks, and when a scenario is held out, 295 of its 849 boundaries have
a label that never occurs in training. ADS-24 met every requirement:
254 runs carry stage labels, of which 132 enter the contract, the
simulated users can be switched on and off, and we downloaded the data
ourselves. The reasons in Table~\ref{tab:ads24-candidates} come from
our own inspection of each candidate's distribution and documentation
on 20 November 2026, not from the cited papers, except where a cited
paper states them.
% 출처(검토일): atlas OT/PROPHET_ENV_CONSENSUS_KD_V3_20261120/EXTERNAL_DATASET_SUITABILITY_V3.md:3;
%   PWNJUTSU_BROWSER_V3_DRAFT.md:3 (ads.md 30번).
% 근거: RECORD:265--275. 판정과 이유는 기록의 표를 그대로 옮겼다.

\begin{table}[htbp]
\centering
\small
\caption{Public datasets considered for the exploratory cross-dataset
check. Reasons record our own inspection on 20 November 2026; sources
without a reference are given in the notes below the table.}
\label{tab:ads24-candidates}
\begin{tabularx}{\textwidth}{@{}>{\raggedright\arraybackslash}p{0.30\textwidth}l>{\raggedright\arraybackslash}X@{}}
\toprule
Candidate & Decision & Reason \\
\midrule
ADS-24/Crate~\cite{nyberg2026cyber} & adopted & 254 runs with stage labels;
  simulated users on and off; downloaded \\
CAM-LDS$^{a}$ & not used & Linux; no benign
  user activity during attacks; 34 runs; with a scenario held out, 295 of
  849 boundaries have a label absent from training \\
Windows multi-source ATT\&CK-labeled data~\cite{niloy2026multi} & not used
  & we could not confirm a public location and license for the raw data
  of the 870 sessions or find action start times \\
APTGen~\cite{takahashi2020aptgen} & not used & distribution site behind
  authentication; no benign background logs \\
DARPA OpTC$^{b}$ & not used & narrative ground truth for one connected
  red-team campaign over three days; we found no way to split it into
  independent runs with stage labels \\
ATLASv2~\cite{riddle2023atlasv2} & not used & ten scenarios; we found no
  per-stage occurrence times in the ground truth \\
PWNJUTSU$^{c}$~\cite{berady2022pwnjutsu} & not used & files behind a login; the dataset paper
  reconstructs the campaign of one participant out of 22, in 18 stages \\
\bottomrule
\end{tabularx}
\par\smallskip
\begin{minipage}{\textwidth}\footnotesize
$^{a}$\url{https://zenodo.org/records/18861762}; checked 20 November 2026.
$^{b}$\url{https://github.com/FiveDirections/OpTC-data}, including its
red-team ground-truth document; checked 20 November 2026.
$^{c}$IEEE DataPort, DOI 10.21227/ngjr-kh92; dataset paper DOI
10.1109/TNSM.2022.3183476; checked 20 November 2026.
\end{minipage}
\end{table}
% 근거: RECORD:267--275.
% 출처(OpTC·PWNJUTSU·CAM-LDS): ads.md 30·31번; atlas V3/OPTC_ATLAS_V3_DRAFT.md:9,11 (저장소, 3일 캠페인);
%   V3/PWNJUTSU_BROWSER_V3_DRAFT.md:3,9,10,17 (검사일, DataPort, P12 한 명 18단계, DOI);
%   V3/PWNJUTSU_BROWSER_V3_DRAFT.md:21-22 (CAM-LDS HEAD 200). 22명은 RECORD:275.
% 표 안 \footnote는 tabularx에서 사라질 수 있어 표 아래 주석으로 옮겼다.
% 7회차 인용 P1: PWNJUTSU 데이터셋 논문(DOI 10.1109/TNSM.2022.3183476)을 임시 key tmp1로 인용했다.
%   tenure가 확인하면 진짜 key로 바꾼다. DOI 표기는 배포본 접근 경로로 남겼다.
% TODO(사람): Wazuh 문서 URL 두 개의 접근일(atlas 근거 없음). 수집에 쓰인 Wazuh 버전은
%   검사한 KTH 기록에 적혀 있지 않다.
